% ===== §12 =====
\section{The Field of Cultivation and the Field of Exercise}
\label{sec:fields}

\noindent
If the capacities of relational understanding cannot, at the depths of intimacy, be rehearsed in the loop of attempt and correction, and yet they can plainly be developed, since some people bring far more of them to their bonds than others, then their development must occur in part elsewhere than in the bond where they are most consequentially exercised. This section draws the distinction on which the practical proposal of the paper rests, between the field in which a capacity is cultivated and the field in which it is exercised, and it holds that in the case of intimacy these come apart.

The confusion the distinction removes is the expectation that intimacy should serve as its own training ground, that one learns to understand one's partner by practicing on one's partner, refining the capacity through the bond itself. For the near-practicable capacities this is partly available, but for the capacities that reach into the wager it is not, and the expectation that it should be is a source of a particular disappointment, the sense that one ought to be getting better at one's bond through sheer repetition and is somehow failing to. The bond is not a practice ground for its own deepest demands. It is the field of exercise, the place where capacities developed in part elsewhere are wagered, and its unrepeatable and slow-returning history is unsuited to the shaping of a capacity through rapid trial.

The field of cultivation, for much of what relational understanding requires, lies in the wider life in which the loop of attempt and correction is available. The recognition of another's world, the entry into it, the construction of a mapping between two structures, the labor of translation, can be developed in every setting where one meets a world not one's own and can see, with reasonable speed, whether one has entered it. They are developed in friendship, in the sustained encounter with a discipline or a tradition foreign to one's own, in the reading of literature that renders a world one does not inhabit, in the practices of teaching and of care where the entry into another's world is the daily work and its success or failure is visible. A person who has developed these capacities across such a life brings them to an intimate bond already formed, and exercises there what was cultivated in fields where the exercise was correctable.

This gives the first half of a distinction that the next part will develop into the model. Some of the capacities of relational understanding are near-practicable. They admit of cultivation in the loop of attempt and correction, in fields outside the intimate bond, and they can be brought to the bond as developed skills. The recognition of the other's world, the entry, the structural mapping, and the translation are of this kind, and the fact that they are cultivable elsewhere is what allows a person to bring understanding to a bond and not only to discover, across the bond's slow history, whether they had it.

The other half of the distinction concerns the capacities that reach furthest into the wager, and these do not admit of cultivation in the same way, for a reason that must be stated precisely. A capacity whose exercise gives feedback only after long delay, into an unrepeatable history, cannot be shaped by the loop of attempt and correction even in principle, since the loop requires the feedback the wager withholds. The throwing of a generative thing into the other's unfolding, the non-intervention that lets a shared dynamic develop, the preservation of a relational experience against a value that may return only much later, are capacities of this kind. They cannot be practiced to reliability, because reliability is precisely what the loop confers and the loop is unavailable. What can be developed, in their place, is a standing disposition and not a skill honed by feedback, a formed way of being toward the other that inclines one to give rightly into the wager without the confirmation a rehearsal would supply. The nature of this disposition, and the sense in which it can be cultivated where a skill cannot, is the matter of the model the next part proposes.
